\section{An example}\label{sec:example}

Take $k=4$ and $\lambda=(10,7,3,1)$. Its first block is $(10,7)$ and its
terminal block is $(3,1)$, so $\bk(\lambda)=7$. Iterating the transfer gives
the following table.

\begin{table}[H]
\centering
\caption{The transfer sequence for $k=4$ and $\lambda=(10,7,3,1)$.}
\label{tab:transfer-example}
\begin{tabular}{ccl}
\toprule
$i$ & $\lambda^{(i)}$ & $\bk(\lambda^{(i)})$ \\
\midrule
0 & $(10,7,3,1)$ & $7$ \\
1 & $(7,6,1)$ & $5$ \\
2 & $(6,3)$ & $4$ \\
3 & $(3,2)$ & $3$ \\
4 & $(2)$ & $2$ \\
\bottomrule
\end{tabular}
\end{table}

Thus $\mathbf{s}(\lambda)=(7,5,4,3,2)$ and
\begin{equation*}
 \Phi_4(\lambda)=(5,5,4,3,2,1,1).
\end{equation*}
This partition is $4$-strict. Its size is $21$, its length is $7$, and
\begin{equation*}
 \ell_4\bigl(\Phi_4(\lambda)\bigr)
 =(5-3)+(2-0)=4=\ell(\lambda),
\end{equation*}
which displays the three statistic equalities in \Cref{thm:main}.
